\section{The general graph realization}\label{general:section}

The blocks of \Cref{offset:section} have two realizations. We first use
one vertex per legal valuation and attach private leaves to identify
its block. This construction has no degree restriction, but gives the
smaller quantitative compiler. Throughout this section \(r\ge3\).
\Cref{reduction:section} proves the game-transfer argument for this
realization. We then construct subcubic graphs with the same affine block
data in \Cref{graph:section}.

\subsection{Valuation vertices and private leaves}

For \(\eps\in\{0,1\}\), create a base vertex \((B,z)\) for every
block \(B\in\cB\) and valuation \(z\in\cZ_B^\eps\). Two base vertices
are adjacent exactly when
\begin{equation}\label{general:adjacency}
 B\ne C,\qquad
 \cR(B)\cap\cR(C)\ne\varnothing,\qquad
 z(\rho)=w(\rho)\quad
 \text{for every }\rho\in\cR(B)\cap\cR(C).
\end{equation}
There are no edges within a block. Agreement is required on
\emph{all} common readouts; disagreement on even one gives a nonedge.

Write \(J=\abs{\cB}\), and let \(N\) be the common number of base
vertices. The base-vertex counts in the two versions agree by
\Cref{lem:block-fibers}.
Number the blocks by \(j(B)\in\{1,\ldots,J\}\) in the same order
in both versions. At each \((B,z)\), attach \(j(B)(N+1)\) private
leaves, adjacent only to that owner. Denote them by
\(\operatorname{leaf}(B,z,\nu)\), where
\(1\le\nu\le j(B)(N+1)\). Call the resulting graphs \(G^0,G^1\),
and let \(V_B^\eps\) contain all base vertices of block \(B\) and
all their leaves. These parts partition the graph.

\begin{proposition}[General realization]\label{prop:general-realization}
For the affine block data at address length \(r\ge3\), the graphs
\(G^0,G^1\) are simple, uncolored, and of the same positive order.
They have \(O(J^3)\) vertices.

Every \(\xi_B\in\cZ_B^1\) determines a fixed bijection
\(\Psi_{B,\xi_B}:V_B^0\to V_B^1\), depending only on \(B,\xi_B\),
by
\[
 (B,z)\longmapsto(B,z+\xi_B),\qquad
 \operatorname{leaf}(B,z,\nu)\longmapsto
 \operatorname{leaf}(B,z+\xi_B,\nu).
\]
For any set \(U\) of blocks, offsets agreeing on all shared readouts
make the union of these maps an isomorphism of the induced subgraphs
on the corresponding parts.

At every history reached by a winning \(K=r+1\)-slot strategy,
a marked base vertex \((B,z)\) is answered by \((B,z')\) with
\(z'\in\cZ_B^1\). If zero base vertices of distinct blocks are
simultaneously marked, their responses agree on every shared
readout, even when all \(K\) slots are occupied.
\end{proposition}

\begin{proof}
The adjacency rule and owner rule introduce no loops or multiple
edges. Each block has nonempty legal fibers of equal cardinality,
so the two graphs have equal positive orders. Each block has at
most \(64\) valuations, a convenient loose constant, hence
\(J\le N\le64J\). The total order is at most
\[
       N+NJ(N+1)=O(J^3).
\]

Translation by a legal affine offset bijects \(\cZ_B^0\) with
\(\cZ_B^1\). All owners in block \(B\) have the same number of
leaves, so the displayed map is a bijection of the entire part.
Within a base block it preserves equality, and there are no base
edges. Between distinct active blocks the offsets agree on every
shared readout; adding equal offsets preserves both agreement and
disagreement on each name, hence the entire condition
\eqref{general:adjacency}. Blocks with no common name stay
nonadjacent.

A leaf is adjacent only to its owner. For a base vertex in the
owner's block the same translation is applied to the base and the
owner, whether or not the owner itself is marked, so their equality
is preserved. A base in a different block remains a nonowner.
Distinct leaves stay distinct because owner translation is
injective and the local index is unchanged. Repeated selections,
different leaves at one owner, and leaves at different owners are
therefore all handled. Leaf--leaf adjacency is absent, and no leaf
is identified with a base vertex. This proves the induced-subgraph
claim on any union of compatible parts.

For recognition, first observe that a marked pair reached under
any winning strategy with at least two slots has equal degrees.
After discarding every other pair, request another slot. The
announced bijection must preserve adjacency to the retained pair
for every choice, and therefore matches its two neighborhoods
bijectively. This is a property of the reached history, not a test
that needs spare slots in that history.

Every leaf has degree one. A base vertex in block number \(j\)
has degree in
\[
       [j(N+1),\,j(N+1)+N-1].
\]
These intervals are disjoint and have lower endpoints greater
than one. Thus every reached base response is a legal valuation
in the same block. Distinct homogeneous zero vertices sharing
a readout are adjacent, so their responses must satisfy
\eqref{general:adjacency} and agree on all shared names.
\end{proof}

The proposition supplies both interfaces needed for the game:
compatible offsets induce blockwise maps on \emph{all} vertices,
including leaves, while zero queries decode to legal affine
valuations with shared-readout agreement. The ring-transfer proof
in \Cref{reduction:section} will use only those interfaces.
